\scanpage{117}
% archive foliation: 116
point est obtenu, pour un $i$ et un $z_i$ fixé, par le choix d'un $z'_i$ satisfaisant
\[
\tag{66}
{z'_i}^2=\underbrace{z_i^3+a_1z_i^2+a_2z_i+a_3}_{(z_i-\zeta_i)(z_i-\zeta_\lambda)(z_i-\zeta_\mu)}.
\]
Choisir, pour $\forall i\in\unsure{J|I}$, un $x_i\in E(k)$ tel que
\[
\tag{67}
2x_i=\varepsilon_i
\]
revient donc à choisir, pour $\forall i$, des quantités
\[
\tag{68a}
\wp(x_i)=z_i
\]
soumises à (65) \inserted{(2 choix possibles)}, et ensuite, la quantité $z'_i$
\[
\tag{68b}
\wp'(x_i)=z'_i,
\]
soumise à (66) --- ce qui fait quatre choix possibles pour $\forall i$, comme il se doit (car l'indétermination pour $x_i$ est dans $ {}_2E(k)$).

Ceci posé la condition d'épinglage de Jacobi relatif d'échelon $4$
\[
\tag{69}
\sum_{i\in\unsure{J|I}}x_i=0
\]
s'exprime par la condition
\[
\tag{70}
\left\{\begin{aligned}
\det\begin{pmatrix}
1&z_1&z'_1\\
1&z_2&z'_2\\
1&z_3&z'_3
\end{pmatrix}&=0\qquad\text{i.e.}\\
(z_1z'_2-z_2z'_1)+(z_2z'_3-z_3z'_2)+(z_3z'_1-z_1z'_3)&=0.
\end{aligned}\right.
\]
Cf prop. \unsure{3)}, qui s'applique, car les $x_i$ sont tous distincts. ---)
